\documentclass[11pt]{article}
\usepackage[margin=1in]{geometry}
\usepackage{graphicx}
\usepackage{amsmath,amssymb}
\usepackage{booktabs}
\usepackage{natbib}
\usepackage{hyperref}
\usepackage{xcolor}

\title{RNAseqCovarImpute: A Multiple Imputation Procedure for \\ Differential Expression Analysis with Missing Covariate Data}
\author{Anonymous}
\date{}

\begin{document}
\maketitle

\begin{abstract}
Missing covariate data is a pervasive and under-addressed problem in observational transcriptomic studies. Standard multiple imputation (MI) procedures handle uncertainty well, but require inclusion of the outcome in the predictor matrix, which is not feasible when the outcome consists of tens of thousands of genes measured on a smaller number of participants. We present RNAseqCovarImpute, an R package that randomly bins genes into small groups, constructs $M$ imputed datasets separately within each bin, fits bin-specific limma-voom models using a globally estimated mean-variance trend, and pools results across imputations with Rubin's rules. Across three simulation studies on 1,044 samples --- with real covariates and real counts, real covariates and seqgendiff-perturbed synthetic counts at 82.5\% and 50\% null-gene rates, and fully synthetic covariates and counts --- RNAseqCovarImpute recovered 95.6--98.5\% of true differentially expressed genes (DEGs) under MCAR missingness (5--55\%) versus 56.6--90.7\% for single imputation (SI), with comparable false discovery rates (FDR) of 2.9--13.8\% and 1.9--14.8\% respectively. Under MAR, complete case (CC) analysis inflated FDR to 11.4--18.5\%. At the 82.5\% null-gene rate, RNAseqCovarImpute identified 96.8--99.1\% of true DEGs under MCAR (SI: 78.9--95.9\%; CC: 58.2--92.9\%). Applied to 1,045 ECHO-PATHWAYS placental samples examining maternal age and controlling for ten covariates, MI identified 382 maternal-age DEGs; CC captured 368 (96\%) and SI 214 (56\%) of these, with 30 DEGs exclusive to MI. RNAseqCovarImpute integrates directly with the limma-voom pipeline and provides a principled handling of missing covariate data for observational RNA-seq studies.
\end{abstract}

\section{Introduction}

Missing covariate data is pervasive in observational studies, and the linear modelling tools that underpin differential expression analysis cannot be fit to records with missing entries. The simplest remedy, complete case (CC) analysis, discards any individual with at least one missing value; this practice reduces statistical power and, depending on the missingness mechanism, can introduce substantial bias into model coefficients \citep{sterne2009multiple,little2019statistical}. A second widely used strategy, single imputation (SI), replaces each missing value with the most likely prediction from a supplementary model. While SI preserves the full sample, it treats imputed values as if they were observed and therefore underestimates standard errors and inflates confidence in downstream estimates \citep{rubin1987multiple,sterne2009multiple}. Multiple imputation (MI) was developed specifically to propagate the uncertainty associated with filling in missing values: $M$ imputed datasets are created, each is analysed separately, and the resulting estimates and standard errors are combined using Rubin's rules \citep{rubin1987multiple,marshall2009combining}. Barnard and Rubin's small-sample adjustment to the pooled degrees of freedom further stabilises inference when the completed fraction is modest \citep{barnard1999smalls}.

Despite decades of methodological development in epidemiology, there has been no concerted effort to identify the most advantageous way of handling missing covariate data in transcriptomic studies. Historically, a large fraction of RNA-sequencing (RNA-seq) work has been performed in in-vitro or in-vivo models that do not suffer from missing covariates, and experimental designs with a small number of controlled variables often avoid the problem entirely. The calculus is changing quickly. As sequencing costs continue to fall, transcriptomic assays are being deployed at scale inside human observational cohorts such as ECHO-PATHWAYS \citep{paulson2022lewinn,paquette2023placental}, where missing covariate data is the rule rather than the exception. Without tailored methods, analysts revert to CC or SI, each with well-understood statistical hazards.

Direct application of MI to RNA-seq is not straightforward. The most effective MI predictor matrices include the outcome variable because doing so preserves the outcome--covariate relationship in the imputed values \citep{vanbuuren2011mice}. For RNA-seq experiments with tens of thousands of genes, however, placing every gene in the predictor matrix would require an equally large number of participants and is typically infeasible. An obvious alternative --- constructing one set of $M$ imputed datasets per gene --- quickly becomes computationally prohibitive. Analogous challenges arose in epigenome-wide association studies (EWAS), where methods based on grouping CpG sites together in the predictor matrix were introduced to recover tractability \citep{mill2013epigenomewide,weinhold2020statistical}. Inspired by these ideas, we developed \emph{RNAseqCovarImpute}, an R package that bins genes into small groups, constructs $M$ imputed datasets separately within each bin, and pools results with Rubin's rules. The package integrates with the limma-voom pipeline \citep{law2014voom,ritchie2015limma}, the de facto standard for bulk RNA-seq differential expression. We assess RNAseqCovarImpute in three simulation studies and a real-world application to placental transcriptomics.

\paragraph{Related work.} The limma-voom framework estimates gene-wise precision weights from the mean-variance trend of log counts per million (log-CPM) and uses those weights inside limma's moderated linear model \citep{law2014voom,ritchie2015limma}. Alternative pipelines such as edgeR \citep{robinson2010edger} and DESeq2 \citep{love2014deseq2} fit negative binomial models directly to integer counts; normalisation by trimmed mean of M-values (TMM) \citep{robinson2010tmm} is standard across all three. Enrichment analyses such as camera \citep{wu2012camera} and GAGE \citep{luo2009gage} operate on gene-level statistics or rankings, which will matter for compatibility with MI approaches. Single-cell RNA-seq has motivated a large body of work on imputing missing or zero counts in the outcome matrix to address dropout \citep{hao2024scrnaseq}, but those methods target sparsity in the outcome rather than missingness in covariates. Generic imputation tools such as mice (chained equations for mixed-type data) \citep{vanbuuren2011mice} and missForest (nonparametric random-forest imputation) \citep{stekhoven2012missforest} operate at the covariate level but were not designed for outcome dimensionalities exceeding sample size. Previous work on MI for EWAS binned CpG sites to make outcome-informed MI tractable \citep{mill2013epigenomewide,weinhold2020statistical}; reported tradeoffs favoured MI for sensitivity but at the cost of higher FDR in some settings. Methylation arrays such as the Illumina MethylationEPIC BeadChip measure around 850,000 sites \citep{pidsley2016critical}, an order of magnitude more than a typical RNA-seq experiment. Finally, we evaluate our method using both real ECHO-PATHWAYS counts \citep{paulson2022lewinn} and synthetic counts generated with seqgendiff \citep{gerard2020data}, which uses binomial thinning of a real count matrix to induce user-specified differential expression coefficients while preserving realistic variance structure; FDR control is performed with the Benjamini--Hochberg procedure \citep{benjamini1995controlling}.

\section{Materials and Methods}

\subsection{The RNAseqCovarImpute procedure}

RNAseqCovarImpute performs multiple imputation and differential expression in six linked steps (Figure~\ref{fig:method}). First, the gene set is randomised and partitioned into bins. By default, each bin contains approximately one gene per ten participants --- a study with 500 participants and 10,000 genes produces 200 bins of 50 genes --- but the ratio is user-configurable. When the gene count is not divisible by the bin size, the procedure creates bins of two adjacent sizes; for 10,001 genes it would emit 199 bins of 50 and one bin of 51. Second, $M$ imputed datasets are generated separately within each bin using the mice R package \citep{vanbuuren2011mice}. The imputation predictor matrix for a given bin consists of all covariates plus the log-CPM values of the genes assigned to that bin. Default mice settings apply: predictive mean matching for continuous variables, logistic regression for binary variables, polytomous regression for unordered categorical variables, and proportional-odds modelling for ordinal variables. The user may specify $M$ (number of imputations) and \texttt{maxit} (number of iterations).

\begin{figure}[t]
  \centering
  \includegraphics[width=0.98\linewidth]{figures/fig1.pdf}
  \caption{Overview of the RNAseqCovarImpute workflow. (A) The raw count matrix (example: 10,000 genes $\times$ 500 samples) is (B) partitioned into small gene bins (default: approximately one gene per ten samples). (C) $M$ imputed datasets are constructed independently within each bin using mice, with the log-CPM values of the bin's genes included in the imputation predictor matrix. (D) Limma-voom with a globally fit mean-variance trend is applied per bin per imputation. (E) Gene-level results are un-binned and stacked across imputations. (F) squeezeVar empirical-Bayes variance shrinkage is applied per imputation. (G) Rubin's rules pool coefficients, standard errors, and Barnard--Rubin adjusted degrees of freedom. (H) The pooled P-values are FDR-adjusted.}
  \label{fig:method}
\end{figure}

Third, voom precision weights are estimated from the mean-variance trend using \emph{all} genes across \emph{all} $M$ imputations rather than separately within each bin. Concretely, lmFit \citep{ritchie2015limma} is run on each bin-imputation pair, the residual standard deviations and average log-counts are extracted, and a robust LOWESS fit across this pooled set defines the global precision weights \citep{law2014voom}. A modified voom function then injects these global weights into bin-specific lmFit calls, which estimate weighted linear models for every gene while respecting the experimental design and covariate adjustment. Fourth, the $M$ sets of bin-level results are un-binned and restacked into $M$ tables of gene-wise model fits. Fifth, squeezeVar is applied separately within each of the $M$ tables to perform empirical-Bayes variance shrinkage with gene-wise prior variances and degrees of freedom. Sixth, Rubin's rules \citep{rubin1987multiple,marshall2009combining} pool coefficients and standard errors across the $M$ tables; the Barnard--Rubin small-sample adjustment to the pooled degrees of freedom is applied \citep{barnard1999smalls}. Wald statistics are computed as pooled coefficients divided by pooled standard errors, two-sided P-values are derived from the $t$-distribution, and P-values are FDR-adjusted following Benjamini--Hochberg \citep{benjamini1995controlling}. Pooling is performed separately for the raw and limma-moderated lmFit results, allowing the user to inspect both.

\subsection{Data and preprocessing}

Real placental RNA-seq and covariate data come from the ECHO prenatal and early childhood pathways to health (ECHO-PATHWAYS) consortium \citep{paulson2022lewinn,paquette2023placental}. Among the 1,503 participants enrolled in the consortium, transcriptomic data are available for 1,083 individuals; we excluded 18 placental abruptions and 20 individuals missing maternal age data, leaving 1,045 for the maternal-age application and 1,044 for the simulation experiments (one additional individual missing prenatal tobacco data was removed for the simulation). Only protein-coding genes, processed pseudogenes, and lncRNAs were retained, and genes with average log-CPM below zero were filtered; this yielded 14,027 genes for the simulation dataset and 14,029 genes for the application. Log-CPM values were TMM-normalised to library size \citep{robinson2010tmm}. Positive prenatal tobacco exposure was coded from a maternal urine cotinine threshold of 200 ng/mL combined with self-reported tobacco use \citep{thyssen2011smoke}.

\subsection{Simulation design}

We compared RNAseqCovarImpute with SI (missForest; \citealp{stekhoven2012missforest}) and CC analysis across three simulation studies. In all three, we first ran limma-voom on the entire sample with complete covariate data (the ``full data'' model) and defined true DEGs as genes with FDR $<0.05$ from that model. We then induced missingness only in covariates using mice::ampute, never in the predictor of interest, at levels of 5--55\% under MCAR and MAR mechanisms, with ten replicate datasets per scenario. CC dropped any individual with at least one missing covariate; SI filled covariates with missForest; MI used RNAseqCovarImpute. The limma-voom pipeline was then applied for each method, and we recorded the percent of true DEGs recovered, the FDR (defined as the percent of false positives among reported significant genes relative to the full-data truth), and the bias of estimated coefficients relative to the full-data coefficients.

Simulation~1 used the real ECHO-PATHWAYS placental covariates and real RNA-seq counts ($N=1{,}044$; 14,027 genes), with maternal age as the predictor of interest and fetal sex, prenatal tobacco exposure, and RNA-seq batch as covariates. Simulation~2 modified the counts with seqgendiff \citep{gerard2020data} to install known signal: binomial thinning adjusted the count matrix so that the coefficients for 2,000 randomly sampled genes matched values drawn from a gamma distribution, with 82.5\% of genes set to null and the remainder drawn from a non-null gamma (matching the approximately 82.5\% non-significant rate in the real data). A sensitivity analysis repeated the exercise with only 50\% of coefficients set to null. Simulation~3 replaced covariates with fully synthetic draws: five normal variables $V_1,\dots,V_5$ with mean 0 and standard deviation 1, with $V_1$ the predictor of interest and $V_2$--$V_5$ the confounders, were generated so that Spearman rank correlations between $V_1$ and $V_2$--$V_5$ targeted 0.1, with no correlation among $V_2$--$V_5$. Counts were thinned as in Simulation~2. The simulation was run at sample sizes of 500, 200, 100, and 50, and sensitivity analyses at $N=200$ varied the correlation between $V_1$ and $V_2$--$V_5$ over 0.2, 0.3, and 0.4.

\subsection{Real-world application}

We applied RNAseqCovarImpute to examine the association of maternal age with the placental transcriptome in the ECHO-PATHWAYS sample ($N=1{,}045$; 14,029 genes) while adjusting for ten covariates: family income (adjusted for region and inflation, USD), maternal race (Black vs. Other), maternal ethnicity (Hispanic/Latino vs. not Hispanic/Latino), maternal education (less than high school vs. high school completion vs. college or technical school vs. graduate/professional degree), study site, maternal alcohol use during pregnancy (yes vs. no), maternal tobacco use during pregnancy (yes vs. no), delivery method (vaginal vs. Cesarean section), fetal sex (male vs. female), and RNA-seq batch. American Indian/Alaska Native, Multiple Race, and Other categories were collapsed with White participants to avoid empty cells in the multivariable model. We ran CC, SI (missForest), and MI (RNAseqCovarImpute) in parallel, with significance declared at FDR-adjusted $P<0.05$.

\section{Results}

\subsection{Simulation 1: real covariates and real counts}

The full-data limma-voom model identified 2,453 true DEGs associated with maternal age while controlling for covariates. When missingness was MCAR, both imputation methods identified more true DEGs than CC, and RNAseqCovarImpute recovered the largest share across all levels of missingness (Figure~\ref{fig:sim1}). The percent of true DEGs detected ranged from 95.6\% to 98.5\% for RNAseqCovarImpute and from 56.6\% to 90.7\% for SI; FDR was comparable across methods (RNAseqCovarImpute 2.9--13.8\%, SI 1.9--14.8\%, CC 5.3--12.8\%). Bias was more tightly centred around zero for RNAseqCovarImpute, both across all 14,027 genes and across the 2,453 true DEGs specifically.

\begin{figure}[t]
  \centering
  \includegraphics[width=0.98\linewidth]{figures/fig2.pdf}
  \caption{Simulation 1: real ECHO-PATHWAYS placental covariates and real RNA-seq counts ($N=1{,}044$; 2,453 full-data true DEGs). Bars show the range of each metric across missingness levels from 5\% to 55\% under MCAR (top) and MAR (bottom). RNAseqCovarImpute recovers 95.6--98.5\% of true DEGs under MCAR and 90.3--97.4\% under MAR, versus 56.6--90.7\% and 55.7--77.6\% for single imputation. FDR is comparable across methods under MCAR but rises for complete case under MAR (CC 11.4--18.5\%).}
  \label{fig:sim1}
\end{figure}

Under MAR, both imputation methods again outperformed CC on true DEG recovery and FDR. RNAseqCovarImpute recovered 90.3--97.4\% of true DEGs versus 55.7--77.6\% for SI. MI and SI had similar FDR (2.0--8.0\% and 1.4--12.2\%), but CC FDR climbed to 11.4--18.5\%. Bias was again more tightly centred around zero for RNAseqCovarImpute than for SI or CC, with the largest discrepancies for the CC results.

\subsection{Simulation 2: real covariates, synthetic counts}

Using seqgendiff to install known differential expression signal with an 82.5\% null-gene rate, RNAseqCovarImpute consistently achieved the highest true DEG recovery (Figure~\ref{fig:sim2}). As an initial diagnostic, we confirmed that the coefficients estimated from the limma-voom pipeline on the thinned count matrix with no missingness closely matched our user-defined coefficients (Supplemental Figure~1). Under MCAR, the percent of true DEGs identified ranged from 96.8\% to 99.1\% for RNAseqCovarImpute, from 78.9\% to 95.9\% for SI, and from 58.2\% to 92.9\% for CC. FDR was similar for MI (1.4--6.1\%) and SI (1.4--7.6\%), while CC ran higher (3.0--8.0\%). Bias was more tightly centred around zero for RNAseqCovarImpute than for SI or CC.

\begin{figure}[t]
  \centering
  \includegraphics[width=0.98\linewidth]{figures/fig3.pdf}
  \caption{Simulation 2: real covariates with seqgendiff synthetic counts at an 82.5\% null-gene rate ($N=1{,}044$). Under MCAR, RNAseqCovarImpute recovers 96.8--99.1\% of true DEGs versus 78.9--95.9\% for SI and 58.2--92.9\% for CC, with comparable FDR (MI 1.4--6.1\%, SI 1.4--7.6\%) and higher CC FDR (3.0--8.0\%). MAR patterns mirror MCAR, with true-DEG recovery 94.0--98.2\% (MI), 77.6--89.6\% (SI), 61.1--93.2\% (CC) and FDR 0.9--3.9\%, 1.2--7.7\%, 6.0--10.9\% respectively.}
  \label{fig:sim2}
\end{figure}

Under MAR with the 82.5\% null rate, the ordering was preserved: true DEG recovery was 94.0--98.2\% for MI, 77.6--89.6\% for SI, and 61.1--93.2\% for CC, with FDR of 0.9--3.9\%, 1.2--7.7\%, and 6.0--10.9\% respectively. The bias patterns paralleled MCAR. In a sensitivity analysis with the null-gene rate reduced from 82.5\% to 50\%, all three methods performed slightly better, but the ordering (MI $>$ SI $>$ CC) persisted across MCAR and MAR for true-DEG recovery, FDR, and coefficient bias.

\subsection{Simulation 3: synthetic covariates and synthetic counts}

With fully synthetic covariates at target Spearman rank correlation 0.1 between $V_1$ and $V_2$--$V_5$, RNAseqCovarImpute was the best performer at identifying true DEGs and controlling FDR across sample sizes of 500, 200, 100, and 50 and across both missingness mechanisms. Coefficient bias was more tightly centred around zero for RNAseqCovarImpute than for SI or CC, both across all genes and restricted to the true-DEG subset. In sensitivity analyses with the $V_1$--$V_{2\text{-}5}$ correlation set to 0.2, 0.3, and 0.4 at $N=200$, RNAseqCovarImpute remained the best performer in general. At a correlation of 0.4, SI identified slightly more true positives than RNAseqCovarImpute, but at the cost of substantially higher FDR, which exceeded 60\% in one scenario.

\subsection{Bias behaviour under MAR and MCAR}

A recurring pattern across Simulations~1--3 is that coefficient bias is most tightly centred around zero for RNAseqCovarImpute, regardless of whether bias is computed across all genes or restricted to the true-DEG subset, and regardless of the missingness mechanism. The ordering MI $<$ SI $<$ CC on absolute bias was preserved across missingness levels from 5\% to 55\%, which matches the expectation that outcome-informed MI should better preserve the outcome--covariate relationship in the imputed values than either SI (which ignores MI uncertainty) or CC (which selectively removes participants). The stability of this ordering across mechanisms is practically relevant because the mechanism is rarely known \emph{a priori} in observational data.

\subsection{Application to maternal age and placental transcriptome}

Among the 1,045 ECHO-PATHWAYS participants included in this analysis, 61 (6\%) were missing data for at least one of the ten covariates. The CC, SI, and MI (RNAseqCovarImpute) analyses identified 575, 214, and 382 significant maternal-age DEGs, respectively (Figure~\ref{fig:app}). CC recovered 368 of the 382 MI-identified DEGs (96\%) and SI recovered 214 (56\%); 30 DEGs were exclusive to MI. Although there were differences, P-value rankings for each method largely preserved the same ordering for the top-ranked genes. The top DEGs included genes with well-characterised roles in inflammation and innate immunity: S100A12 and S100A8 are pro-inflammatory calcium-, zinc-, and copper-binding proteins; CXCL8 (IL-8) and IL1R2 are pro-inflammatory cytokine and cytokine-receptor genes; SAA1 and CASC19 are expressed in response to inflammation; and LILRA5, a leukocyte receptor gene, has been implicated in triggering innate immune responses.

\begin{figure}[t]
  \centering
  \includegraphics[width=0.98\linewidth]{figures/fig4.pdf}
  \caption{Application to maternal age and the placental transcriptome ($N=1{,}045$; 6\% of individuals missing at least one of ten covariates). The complete case, single imputation, and multiple imputation analyses identify 575, 214, and 382 differentially expressed genes, respectively (left). Of the 382 MI-identified DEGs, CC captures 368 (96\%) and SI captures 214 (56\%); 30 DEGs are MI-exclusive (right).}
  \label{fig:app}
\end{figure}

\subsection{Summary across simulations}

Aggregating across the three simulation studies, two robust patterns emerge. First, RNAseqCovarImpute dominated SI and CC on true DEG recovery by a margin that widened as missingness increased: at the highest missingness level tested (55\%), the gap between MI and SI was approximately 20--40 percentage points on true DEG recovery, while the gap was smaller (single digits) at 5\% missingness. Second, the FDR difference between MI and SI was modest throughout, while CC showed larger FDR inflation under MAR --- a mechanism known to bias CC coefficients --- than under MCAR. Together these patterns confirm that the performance advantage of RNAseqCovarImpute is driven primarily by sensitivity gains at moderate-to-high missingness and by its robustness to MAR violations of the stricter MCAR assumption.

\section{Discussion}

Our simulations and real-data application together establish that a multiple imputation procedure that deliberately includes outcome genes in the imputation predictor matrix --- made tractable by the gene-binning strategy introduced in RNAseqCovarImpute --- delivers a consistent performance advantage over complete case and single imputation approaches for RNA-seq differential expression analysis. Across three simulation studies and multiple missingness mechanisms, RNAseqCovarImpute recovered a larger fraction of true DEGs than either alternative while controlling FDR at levels comparable to single imputation and limiting coefficient bias more tightly than any other method tested. The magnitude of the advantage was largest when CC was the comparator: under MAR in Simulation~1, CC FDR climbed to 11.4--18.5\% while MI held FDR at 2.0--8.0\%. These numbers should be interpreted in the context of observational RNA-seq studies, where the ten-covariate model in our placental application is not unusual and 6\% of participants were missing at least one covariate --- a scenario that would cause CC analyses to silently discard dozens to hundreds of samples across realistic cohorts.

A natural point of comparison is prior work on multiple imputation for epigenome-wide association studies, where binning CpG sites in the predictor matrix made outcome-informed MI tractable \citep{mill2013epigenomewide,weinhold2020statistical}. In those studies, MI often uncovered more true positives than CC but at the cost of higher FDR. Our results tell a different story for RNA-seq: RNAseqCovarImpute recovered more true DEGs \emph{without} an FDR tradeoff. We attribute this difference partly to the lower dimensionality of the transcriptome relative to the methylome (tens of thousands of genes versus approximately 850{,}000 CpG sites on the Illumina MethylationEPIC BeadChip \citep{pidsley2016critical}), which permits smaller numbers of bins and therefore better-behaved imputation models, and partly to the specific choice of using limma-voom's globally fit mean-variance trend across all genes in all imputations when computing precision weights. This design keeps the voom weights coherent even though individual lmFit calls operate on narrow bins.

Although imputation methods have been extensively developed to address sparsity in single-cell RNA-seq --- dropout-style missingness in the outcome matrix \citep{hao2024scrnaseq} --- very little attention has been paid to imputation of missing covariate data in bulk RNA-seq, likely because the field has historically focused on controlled in-vitro and in-vivo experiments with few and easily measured covariates. The calculus is changing rapidly. Observational human cohorts such as ECHO-PATHWAYS \citep{paulson2022lewinn,paquette2023placental} now routinely generate transcriptomes at sample sizes of 1,000 or more, with complex covariate structures inherited from the epidemiological study design. Methods for the treatment of missing data are well-established in this community \citep{sterne2009multiple,little2019statistical,marshall2009combining}, with MI emerging as a preferred tool. Yet the transcriptomics literature has not kept pace: when missing covariates are addressed at all, CC dominates, with SI as a distant second. RNAseqCovarImpute lets observational transcriptomic studies adopt MI without abandoning their differential expression pipeline of choice, provided that pipeline is limma-voom.

The maternal-age application illustrates why the choice of missing-data method matters in practice. The three methods identified overlapping but distinct DEG lists, with 30 genes exclusive to MI and substantial numerical differences in how many DEGs each method detected (575 for CC, 214 for SI, 382 for MI). Biological interpretation of the top-ranked genes converged on inflammation and innate immunity --- S100A12, S100A8, CXCL8 (IL-8), IL1R2, SAA1, CASC19, LILRA5 --- a pattern that would be recognised regardless of the missing-data method chosen, but the exact DEG set is non-trivially different. Beyond gene-list differences, any missing-data strategy is preferable to ignoring confounding altogether. Prior work has shown that unadjusted analyses can reverse sign once confounders are introduced: positive associations of young maternal age with child ADHD in unadjusted models are eliminated or even reversed after adjustment for prenatal tobacco and other mediators. Failing to control for these variables can therefore lead to erroneous clinical recommendations --- e.g. targeting ``ideal age for childbearing'' rather than addressing tobacco exposures that may mediate the association.

The main limitation of RNAseqCovarImpute is downstream compatibility. The method produces a single pooled list of genes with $t$-statistics, log fold changes, and P-values, which is compatible with gene-set enrichment approaches that consume rankings or gene-level statistics --- camera \citep{wu2012camera} and GAGE \citep{luo2009gage} work out of the box --- but is not directly compatible with enrichment methods that require a sample-level expression matrix. The package was designed around limma-voom \citep{law2014voom,ritchie2015limma}; users who prefer edgeR \citep{robinson2010edger} or DESeq2 \citep{love2014deseq2} can still use the binning and imputation steps separately, export the imputed datasets, and plug them into alternative pipelines, but out-of-the-box compatibility is restricted. Future work should extend the pooling machinery to the negative-binomial model class and to matrix-level enrichment methods, which will likely require developing new pooling rules for dispersion parameters and for sample-level bootstraps. Another open question concerns runtime scaling: at the default one-gene-per-ten-sample bin size, the method fits thousands of bin-specific imputation models, and although each is small, their aggregate cost can be non-trivial at $M=20$--$50$ imputations. Parallelisation across bins is an obvious remedy and is easy to implement in the package.

Several practical considerations deserve emphasis. The default bin size of approximately one gene per ten participants is a heuristic; in practice, users with larger cohorts may safely use larger bins (e.g. one gene per five participants) without sacrificing the numerical stability of the bin-level imputation models, while very small cohorts (e.g. $N \leq 100$) may benefit from correspondingly smaller bins. The number of imputations $M$ should scale with the expected fraction of missing information; values of $M$ between 20 and 50 are usually adequate for FMI below 50\%, and larger values add diminishing returns. The choice of imputation method for each covariate type is left to mice defaults, which are sensible for most epidemiologic datasets; users with exotic covariate distributions (e.g. zero-inflated continuous variables) may supply their own method specifications. Finally, users should always inspect convergence diagnostics of the mice chains before trusting downstream inference, exactly as they would in any non-transcriptomic MI analysis.

Another practical consideration concerns the interpretation of the 30 MI-exclusive DEGs identified in the placental application. These genes are significant in the pooled MI analysis but not under CC or SI. Three non-exclusive explanations are possible: (i) they are true positives that CC missed due to sample discarding and that SI missed due to over-confident standard errors; (ii) they are Type~I errors driven by between-imputation variability; or (iii) they reflect genuine differences in the maternal-age--gene-expression relationship that are partially obscured by missing confounders. Discriminating among these requires biological follow-up --- replication in independent cohorts, mechanistic experiments, and pathway-level analyses --- but the methodological fact that the method identifies them at all, and holds FDR comparable to SI, is consistent with genuine recovery of signal rather than inflation of noise.

\section{Conclusions}

As sequencing costs continue to decrease, RNA-seq studies are shifting from small-scale in-vitro and in-vivo experiments to large-scale human observational cohorts where missing covariate data is prevalent and consequential. Multiple imputation is the statistical method of choice in epidemiology, but its direct application to transcriptomic outcomes has been blocked by the high dimensionality of the outcome matrix. RNAseqCovarImpute bridges this gap by binning genes, imputing within bins, and pooling with Rubin's rules. In three simulation studies at $N=1{,}044$ and in a real application at $N=1{,}045$, our method recovered more true DEGs than single imputation or complete case analysis, controlled FDR at comparable or better rates, and reduced coefficient bias. The R package integrates directly with limma-voom and exposes separate binning, imputation, model-fitting, and pooling functions so that users can build custom pipelines. Taken together, these results make a concrete case for adopting multiple imputation as the default strategy for handling missing covariate data in observational RNA-seq studies. More broadly, they illustrate a general design principle: when the outcome is too high-dimensional to include directly in an imputation predictor matrix, binning the outcome into groups small enough for the per-bin predictor matrix to remain well-conditioned can recover the benefits of outcome-informed MI without sacrificing tractability. The same principle may extend to proteomics, metabolomics, and other high-dimensional observational assays where missing covariates remain an under-treated problem.

\section*{Availability}

RNAseqCovarImpute is distributed as an R package that integrates with the limma-voom differential expression pipeline. The package exposes separate functions for gene binning, imputation within bins, limma-voom fitting across bins and imputations, and Rubin-rule pooling, which lets advanced users swap in alternative differential expression pipelines such as edgeR or DESeq2 at the cost of writing their own pooling glue. Example scripts reproducing the simulation and application results are bundled with the package so that prospective users can evaluate the method on their own data before production use.

\bibliographystyle{plainnat}
\bibliography{references}

\end{document}
